\label{sec:phase-k}

A natural follow-up to the trade-off characterisation: \emph{can we
use both heads at inference time to produce a sample that is both
faster than pure AR and better than pure diff?} Phase~J showed the
two heads each have unique structural advantages (diff-fill speedup,
FIM capability, revision NLL improvement). Phase~K tests a
\textbf{per-token cross-mode routing} recipe that uses both:

\begin{enumerate}
  \setlength\itemsep{0pt}
  \item AR head generates a prefix of $k_\text{ar}$ tokens.
  \item Diff head fills the next $k_\text{diff}$ tokens in $N$
        denoising steps (parallel, one forward per step).
  \item Per-position commit-time confidence is captured during the
        diff fill: the top-1 probability at the step each position
        gets unmasked.
  \item Bottom-$K\%$ of positions by confidence are flagged.
  \item AR head refills only the flagged positions
        (one teacher-forced forward each).
\end{enumerate}

This is novel at $<\!1$B text-LM scale. \citet{spasov2025deer} (DEER,
2512.15176) drafts with diffusion and verifies with AR but at
\emph{block} granularity. \citet{im2025idlm} (I-DLM, 2604.11035)
and \citet{corrective_dlm_2025} (Corrective DLMs, 2512.15596) do
per-token correction but stay within the diffusion head. Per-token
cross-mode routing — diff drafts in parallel, AR refills the worst
positions identified by confidence ranking — has not been published
at this scale.

\paragraph{The signal: what works and what doesn't.}
We tried four routing signals on F10 final:

\begin{center}
\begin{tabular}{lll}
\toprule
Signal & Result & Why \\
\midrule
Diff single-mask placed-token prob & useless (mean $\sim 0.01$) & marginal $P$ at 50 k vocab too diffuse \\
Diff single-mask top-1 agreement & useless (97\,\% disagree) & calibration mismatch (joint vs single mask) \\
AR-verifier teacher-forced prob & useless (mean $\sim 0.003$) & AR and diff have divergent distributions despite shared backbone \\
\textbf{Commit-time conf, percentile rank} & \textbf{works} & relative ordering picks worst placements even if absolute level is low \\
\bottomrule
\end{tabular}
\end{center}

The percentile-based signal is the key insight. \textbf{All diff
confidences are absolutely low} at 305\,M (mean $\sim 0.03$), so any
absolute threshold flags 100\,\% of positions. But the \emph{relative
ordering} encodes information about which placements are worst
within a specific generation; flagging the bottom-$K\%$ produces a
sensible refill rate that AR can productively fix.

\paragraph{Results on F10 final, $N=50$ held-out GSM8K problems.}

\begin{center}
\begin{tabular}{lrrrr}
\toprule
Config & Refill\,\% & Mean NLL & Wall (s) & Loop \\
\midrule
\code{pct10} & $9.4$\,\% & $12.16$ & $4.1$ & $6$\,\% \\
\code{pct25} & $25$\,\% & $11.89$ & $4.3$ & $6$\,\% \\
\textbf{\code{pct50}} & $\mathbf{50}$\,\% & $\mathbf{11.66}$ & $\mathbf{5.0}$ & $12$\,\% \\
\code{pct75} & $75$\,\% & $12.00$ & $5.6$ & $16$\,\% \\
\midrule
\multicolumn{5}{l}{\emph{Phase J.2 baselines (same F10, same prompts)}} \\
\midrule
\code{ar\_128} (pure AR) & --- & $12.56$ & $\sim 6.4$ & --- \\
\code{ms\_96\_32} & --- & $11.82$ & --- & --- \\
\code{ms\_64\_32} & --- & $11.35$ & --- & --- \\
\bottomrule
\end{tabular}
\end{center}

\paragraph{Headline.}
\textbf{At the same generation length (128 tokens), \code{pct50}
beats pure AR by $-0.90$ NLL/token AND is $1.3\times$ faster.}
Composite-specific recipe: pure AR fundamentally cannot do this
(no diff head to draft with). The recipe does not beat
\code{ms\_64\_32} ($11.35$) per-token, but \code{ms\_64\_32}
generates only $64$ tokens, not $128$; the two configurations target
different output lengths.

\paragraph{What this proves.}
The user's intuition — \emph{diff fills the easy structure in
parallel, AR fixes the suspect tokens} — holds at $305$\,M scale,
provided the routing signal is percentile-ranked commit-time
confidence rather than an absolute threshold. The recipe yields a
Pareto improvement over pure AR (better NLL \emph{and} faster
wall-clock) that no single-head architecture can replicate.

\paragraph{What this does not prove.}
We did not measure GSM8K-dev free-run accuracy here; Phase~I
already established that downstream routing does not lift accuracy
at our scale, and Phase~K targets a different metric (per-token
NLL on gold continuation, plus wall-clock). Whether the
$-0.90$ NLL improvement converts to downstream task wins remains
open, as does generalisation to larger scales and non-math
substrates.
